\subsection{ニュートン法型}\label{節：ニュートン法型}

\noindent\textbf{【前提と学ぶこと】}\par
微分係数$f'(a)$が接線の傾きを表すことを使う.
微分が未習なら,接線と$x$軸の交点から次の近似値を作る図形的な発想と,得られた漸化式の変形を先に読んでよい.
閉じた一般項を求める計算と,近似がどれほど速く改善するかを調べることは別の目標である.
後者は\sectionref*{節：差分から始まる離散的な微積分}で見直す.

少々離れたところから話を始める.
ある関数$f(x)$について,\,$f(a)$が既知の時,その近傍の点$x=a+h$で
\[
  f(a+h) \approx f(a)+f'(a)h
\]
が成り立つ.
つまり,\,$f(a)$の値を元に,\,$f(a+h)$を近似したのである.
方程式
\[
  f(x)=x^2-2=0
\]
の正の解 $\sqrt{2}$ を求めたいとする.いまの近似値を $x=a_n$ とし,
点 $(a_n,f(a_n))$ における接線を引く.その接線の傾きは
$f'(a_n)=2a_n$ だから,接線の方程式は
\[
  y=f(a_n)+2a_n(x-a_n)
\]
である.\,$y=0$ と置いて $x$ 軸との交点を求めると,
\begin{align*}
  a_{n+1}
    &=a_n-\frac{f(a_n)}{f'(a_n)}\\
    &=a_n-\frac{a_n^2-2}{2a_n}\\
    &=\frac{a_n^2+2}{2a_n}
\end{align*}
これが
\[
  a_{n+1}=\frac{a_n^2+2}{2a_n}
\]
という,ニュートン法を$f(x)=x^2-2$に適用したときの漸化式である.

例えば $a_1=2$ とすると,
\[
  a_2=\frac{3}{2},
  \quad
  a_3=\frac{17}{12},
  \quad
  a_4=\frac{577}{408}
\]
となり,これらは $\sqrt{2}$ に急速に近づく.この速さは,誤差を
$e_n=a_n-\sqrt{2}$ と置くと見える.
\begin{align*}
  e_{n+1}
    &=\frac{a_n^2+2}{2a_n}-\sqrt{2}\\
    &=\frac{(a_n-\sqrt{2})^2}{2a_n}\\
    &=\frac{e_n^2}{2a_n}
\end{align*}
誤差がほぼ二乗されるため,正しい近似に近づくと桁数が一気に増える.

実は,この漸化式は閉じた式で表すことができる.
\[
  a_{n+1}-\sqrt{2}=\frac{(a_n-\sqrt{2})^2}{2a_n},
  \qquad
  a_{n+1}+\sqrt{2}=\frac{(a_n+\sqrt{2})^2}{2a_n}
\]
であることに注目すると,両辺の比を取ることで$2a_n$が消え,
\[
  \frac{a_{n+1}-\sqrt{2}}{a_{n+1}+\sqrt{2}}
  =\left(\frac{a_n-\sqrt{2}}{a_n+\sqrt{2}}\right)^{2}
\]
という関係が得られる.つまり
\[
  b_n=\frac{a_n-\sqrt{2}}{a_n+\sqrt{2}}
\]
と置くと,\,$b_{n+1}=b_n^{2}$という,対数型(\sectionref{節：対数型})と同じ構造の漸化式になる.先ほど見た誤差の2乗のふるまいは,この$b_{n+1}=b_n^2$そのものだったのである.

\,$a_1=2$のとき,\,$b_1=\dfrac{2-\sqrt2}{2+\sqrt2}=3-2\sqrt2=\dfrac1{3+2\sqrt2}$であるから,\,$u=3+2\sqrt2$とおくと$b_n=u^{-2^{n-1}}$である.これを$b_n=\dfrac{a_n-\sqrt2}{a_n+\sqrt2}$について$a_n$で解けば,
\[
  a_n=\sqrt{2}\,\frac{u^{2^{n-1}}+1}{u^{2^{n-1}}-1},\qquad u=3+2\sqrt2
\]
という閉じた式が得られる.

このように,今回の$f(x)=x^2-2$に対するニュートン法型漸化式では,特別な変形によって一般項を閉じた式で表すことができる.
しかし,一般のニュートン法型漸化式が同じように閉じた式で表せるとは限らない.
閉じた式が得られなくても,収束先や誤差の減り方を調べることには十分な価値がある.
この解法の主眼は閉じた式を得ることだけにあるのではない.
「接線を使って,次の近似値をどう設計するか」を考え,その設計が
どのような速さで誤差を減らすかを調べることこそが本質である.付録では,この考えを
漸化式の収束,固定点,誤差評価,微分方程式の数値解法へ広げる.
